\section{Paper Summary}

The paper presents \textit{NeuroPathX}, an explainable deep learning framework designed to model interactions between structural MRI features and genetic variation aggregated at the biological pathway level. The methodological problem addressed is multimodal imaging-genetics integration under two competing constraints: capturing complex nonlinear associations while preserving interpretability. This problem is medically relevant because neurological disorders such as Autism Spectrum Disorder (ASD) and Alzheimer's Disease (AD) are heterogeneous and highly polygenic, meaning that disease risk emerges from distributed genetic effects interacting with brain structure.

\paragraph{Positioning with respect to the state of the art.}
Existing imaging-genetics approaches either rely on linear multivariate models, which cannot capture higher-order interactions, or on deep learning methods that operate as black boxes. Latent fusion architectures combine compressed modality-specific embeddings but may obscure direct region-gene relationships, while early-fusion cross-attention models lack constraints ensuring stability and interpretability of learned interactions. NeuroPathX addresses these limitations by combining biologically structured inputs with attention-level regularization. The present paper aims to overcome these limitations by structuring genetic inputs biologically and regularizing the attention mechanism itself.

\paragraph{Main contributions.}
The paper introduces three main novelties. First, it aggregates SNPs into KEGG biological pathways using GWAS effect sizes, shifting interpretation from individual variants to functional systems. Second, it proposes a pathway-guided cross-attention mechanism that models direct interactions between pathway-level genetic representations and ROI-level imaging features. Third, it introduces two attention-level regularization terms: a sparsity loss encouraging selective interactions and a pathway similarity loss promoting cohort-level consistency between patients and controls. The framework is evaluated on two distinct neurological disorders, strengthening claims of generalizability.

\paragraph{Methodology.}
Genetic preprocessing follows a polygenic risk score pipeline: disease-associated SNPs are selected from GWAS studies, mapped to genes, and aggregated into KEGG pathways. Structural MRI features (volume, surface area, cortical thickness, and thickness variability) are extracted from predefined ROIs using FreeSurfer. The core model uses cross-attention in which pathway encodings act as queries over ROI features, enabling explicit modeling of pathway–brain interactions. Unlike standard transformers, softmax is replaced by ReLU in the attention layer to allow suppression of irrelevant pathways. The final loss combines weighted cross-entropy with a KL-based sparsity constraint on attention weights and an $\ell_1$ pathway similarity loss computed between patient and control attention summaries.

\paragraph{Validation and results.}
The method is evaluated using 10-fold cross-validation on the ACE dataset for ASD ($N=165$) and the ADNI dataset for AD ($N=438$). Performance metrics include accuracy, sensitivity, specificity, and AUC. NeuroPathX achieves the highest accuracy and specificity on both datasets and outperforms ablated variants without the proposed regularization terms, suggesting that attention constraints improve robustness. Interpretability is assessed by averaging attention matrices across folds to identify dominant pathway–ROI associations, which are qualitatively reported to align with known neurobiological mechanisms in ASD and AD.
